\documentclass[11pt]{article}
\usepackage[margin=1in]{geometry}
\usepackage{amsmath,amssymb}
\usepackage[hidelinks]{hyperref}
\setlength{\parindent}{0pt}
\setlength{\parskip}{0.5em}
\title{Swarm Coordination Under Changing Conditions\\Research Contract Draft}
\author{}
\date{Version 0.5 \quad 24 November 2026}
\begin{document}
\maketitle

\section{Governing question}
\textbf{How can groups of autonomous vehicles keep doing joint work while the conditions for coordination change?} Changes can affect what vehicles observe, what they can do, who may commit shared resources, and what work or safety margin their operators require. The object of study is the transition during joint work, including what remains of earlier commitments and what controllers can know while the change is unfolding.

Fault-tolerant coordination is the first research track. It provides a setting in which to study changes that propagate through commitments and shared resources, including faults that do not involve a deceptive agent. The cascade obligations below define the first integrated return. The other questions remain part of the programme. Legitimate revisions to goals, authority or tolerated risk have their own rules and are not classified as faults merely because they invalidate an earlier plan.

The programme seeks reusable constructions, proofs and limits across these families. That goal does not depend on demonstrating that an ordinary proof cannot establish the same conclusion. Claims of scientific novelty and claims of reduced proof maintenance require different evidence.

\section{Common model and evidence}
Use explicit shared resources, physical state, authority, and local observation and action histories. Name the admissible changes independently of whether a controller succeeds. For each operating envelope, state both a safety invariant and the service owed: deadlines, recurring throughput, or explicit conditions for ending service. Distinguish an issued or accepted commitment from its execution, receipt, discharge, or authorized cancellation. Two histories that end in the same named mode need not leave the same commitments. Memory used to audit those histories need not be readable by a vehicle. When evaluating a substitution, keep the service and observation interfaces fixed unless the comparison charges for changing them.

For anticipated changes on a vehicle, synthesize or design controllers before deployment. One installed policy family must work across locally indistinguishable histories; a separate winner for each true scenario is not a realizability result. An online computation is admissible only where its completion is not silently assumed on the safety-critical path. Finite enumeration and minimal sets of controllers permitted to change are reference calculations, not governing objectives. No earlier private vehicle research is an input.

\section{First track: fault-tolerant coordination}
\label{sec:faults}
Use the \href{https://github.com/eternal-november/eos-research-programmes/blob/e1348964f73704212d2d8b48d5982a5ed600a2ff/projects/byzantine-control/FAULT-COVERAGE.md}{Byzantine Control fault inventory} to construct a model-specific fault set. Classify faults by the operation they change: evidence production or interpretation; transport, clocks or retained messages; decisions and authority; resource supply, effect or actuation; and detection, fallback or recovery. Track common causes and delayed effects across these groups. For each group give a mechanism on the chosen system or a reason it is out of scope. The inventory is a search aid, not a requirement to implement every listed fault. Separate model error from an event admitted by the model. Untrusted or Byzantine members are optional; legal withdrawal and authorization changes retain the distinctions in C8 and C12.

\textbf{Fault coverage.} Report the envelope under which the claims hold, including onset, duration, correlation, observation delay and what survives the incident. Select interacting faults by the commitments and shared resources they can affect, rather than by the number of patterns that can be enumerated. Include an ordinary accidental or environmental case and a common-cause or recovery-induced case in the case register. Only one needs to be the primary C11 witness; a held-out case must exercise a different mechanism as specified in the companion. State service beyond the admitted envelope separately from safety that can be maintained there. A whole-system stop cannot discharge a service obligation.

\section{Specialized governing questions and proposed obligations}

\subsection{C1a. Communication degradation during joint work}
\textbf{Question.} Which forms of joint work survive a change in how vehicles can communicate, and what must change before they can continue?

\textbf{Obligation.} Classify regimes by evidence actually usable at a decision: two-way, intermittent, asymmetric, partitioned or absent exchange, with timing, identity and direction. Include non-RF exchange such as optical, acoustic and docked messaging, intentional signals, persistent marks and mechanical interlocks where the setting supports them. Distinguish passive sensing and prearranged schedules from live messages. Derive an implementable switching rule and a bound on where it works, or prove a realizability limit for a family with live commitments and staggered detection. Vary the history at the cut while holding the destination regime fixed, and vary modality while holding a required commitment fixed. Include a no-usable-shared-evidence case. Compare \href{https://arxiv.org/abs/1603.02023}{timed supervisor localization} and \href{https://arxiv.org/abs/1609.07038}{intermittent communication and route planning}.

\subsection{C1b. Coordinating through local perception}
\textbf{Question.} What coordinated motion or joint service can a group sustain from local perception of other agents, without message exchange?

\textbf{Obligation.} Compare direct bearing or optical motion cues with detection and tracking on a common simulated perception budget. Vary occlusion, field of view, sampling and delay, neighbor identity continuity, and a change in who can see whom during motion. Give a collision and service claim with an observable failure boundary, distinguishing cohesive motion from a guarantee that another agent will honor a commitment. An estimator fed exact identities or full state is a reference, not the local policy. \href{https://pmc.ncbi.nlm.nih.gov/articles/PMC2234121/}{Starling observations} motivate local interaction without establishing vision exclusivity. \href{https://arxiv.org/abs/2012.01245}{Vision-based drone flocking} supplies a detector--tracker--controller baseline. The detector choice is not an obligation.

\subsection{C2. Recurring splits and rejoins}
\textbf{Question.} Can groups repeatedly change who coordinates with whom while earlier commitments still constrain them?

\textbf{Obligation.} Construct a rule for composing successive and overlapping changes, or delimit the class in which such a rule is possible. Handle delayed detection and stale messages after rejoin. Distinguish rebracketing a common event from exchanging the order of two genuine events. Prove which obligations can be retired once they settle, and derive a bound on retained coordination state in terms of live commitments, timing or backlog under a declared recurrence envelope. If a proposed finite account fails, isolate the missing information by a parameterized lower bound or a sharp counterexample with a positive repair. A safe one-time handoff, a common reset or a new simulation state for every epoch is insufficient. Compare with supervisory recovery and the pinned project's recurrent assembly before claiming this is new.

\subsection{C3. Distinct roles and changes of role}
\textbf{Question.} How can heterogeneous members exchange or replace roles as conditions change without duplicating authority, breaking obligations, or silently suspending service?

\textbf{Obligation.} Model a role by the capabilities, observations, permissions and duties it requires. Derive when an available candidate can serve, when authority can pass, and which pending duties must stay with the previous bearer, transfer with evidence, or be discharged first. Cover a role needed continuously, such as relay coverage, and an exclusive role, such as controller of a shared route. An overlap that preserves coverage may be wrong for exclusivity. Require one planned transfer, one partially observed loss, and a second role change before an earlier duty retires. Hold the assignment and candidate capabilities fixed in a matched pair where the outstanding work changes which handoff is feasible. A matching or task-allocation output is not a handoff proof. Compare with \href{https://robotics.usc.edu/publications/347/}{multi-robot task allocation} and \href{https://arxiv.org/abs/2105.05586}{resilient energy-aware task allocation}; characterize or implement an additional guarantee only where those baselines do not already supply it.

\subsection{C4. Interoperability and replacement of partners}
\textbf{Question.} What must independently implemented cooperative groups disclose so they can join, leave, or change partners without exposing or re-verifying all their internal controllers?

\textbf{Obligation.} Derive interface conditions from separately supplied implementations and shared physical constraints. Prove that substitution preserves the consumed information, safety and timed-service judgments through an actual transition with pending old obligations. Distinguish sufficient compatibility from exactness: if claiming exactness, give the converse lifting argument or prove a boundary. Test one replacement with a different sensor, clock or fallback discipline and one that preserves standalone behavior but fails at the joint boundary. Record what checks are reused, what must be done again, and the cost of building the interface. A summary that embeds a solved global system is disallowed. Use \href{https://web.cs.wpi.edu/~heineman/html/teaching_/CS562/p109-de_alfaro.pdf}{interface automata} and \href{https://arxiv.org/abs/2205.12085}{information-flow guided synthesis} as comparators, rather than presenting interface-based reasoning as a new idea.

\subsection{C5. Actions that create evidence}
\textbf{Question.} When should a swarm spend time or maneuvering freedom to make an uncertainty resolvable, rather than act on what it already knows?

\textbf{Obligation.} Put a probe, repositioning, intentional visible action or rendezvous in the controller's action set, with a real effect on sensing, clearance, resource use and later decisions. In a common environment, compare doing without the information, buying it, and reserving a conservative fallback. Derive a threshold, bounds or a policy comparison in terms of decision slack and the evidence actually gained. Include complementary local observations that do not become common merely because every vehicle obtains one. Study \href{https://pmc.ncbi.nlm.nih.gov/articles/PMC12738390/}{planned synchronization with active observations} and \href{https://papers.neurips.cc/paper_files/paper/2020/hash/9db6faeef387dc789777227a8bed4d52-Abstract.html}{decentralized active perception}. The proposed increment concerns evidence used to alter or discharge a joint commitment, not active sensing by itself.

\subsection{C6. Joint risk under degraded information}
\textbf{Question.} When the evidence available to groups changes, what joint risk guarantee can they still justify and what must they change to sustain it?

\textbf{Obligation.} Fix the loss, horizon and admissible joint laws first. Specify whether the guarantee is conditional on evidence at decision time, ex ante over a mission, or uniform over admitted laws. Characterize a class in which locally maintained certificates combine into a joint one across repeated changes; identify precisely which dependence or information loss invalidates the rule and what evidence, resources or policy changes repair it. For a small test, two fair independent backup-failure bits \(X,Y\) and a fair missing-report bit \(L\) have the same pairwise marginals whether \(L\) is independent or \(L=1\) exactly when \(X=Y\). Given \(L=1\), the failure probability of both backups is \(1/4\) in the first case and \(1/2\) in the second. This is a regression, not a frontier claim. Compare existing \href{https://ieeexplore.ieee.org/document/8744598}{risk bounds under packet loss} and \href{https://doi.org/10.1007/s10514-025-10231-6}{risk-constrained rendezvous under intermittent connectivity}. A new result must say more than that a risky policy can use a more conservative threshold.

\subsection{C7. Predictable changes, avoidance and preparation}
\textbf{Question.} When is a forecast best used to avoid exposure, and when should the swarm change its operating arrangement in advance so it can stay exposed safely?

\textbf{Obligation.} Compare staying put with a reactive fallback, always preparing, forecast-directed preparation, and avoidance. Give matched methods the same forecast and its error envelope; vary who receives it, its lead time, missed and false warnings, and the cost and reversibility of precautions. Where action changes exposure, use a causal model to evaluate different routes, not the realized outcome from the route that happened to be taken. Seek a regime comparison or a bound that shows when preparations interact through capacity or inherited obligations. This is the optional allostatic branch: anticipation can change admission, reserves or formation while ordinary feedback continues within the chosen arrangement. Compare with \href{https://arxiv.org/abs/2103.05526}{distributed predictive control with information-age forecasts}; anticipation by itself is not an advance.

\subsection{C8. Multiple owners, withdrawal and limited trust}
\textbf{Question.} What joint service remains possible when independent owners may decline new work, withdraw resources, or supply evidence of differing trustworthiness?

\textbf{Obligation.} Separate legitimate refusal, accidental failure, and arbitrary deviation. Distinguish what an owner is permitted to do from what its vehicles can physically do. Prove which inherited duties can be discharged within the permitted withdrawal model and when an independently effective enforcement or reserve mechanism is necessary. Include a case where revoking a logical permission does not end an already initiated physical use. Import a Byzantine Control result only after aligning its authority, fault, timing and service assumptions. The \href{https://github.com/eternal-november/eos-research-programmes/blob/e1348964f73704212d2d8b48d5982a5ed600a2ff/projects/byzantine-control/research/PROGRAM-CONTRACT.md#L13-L37}{Byzantine Control programme} is a starting point, not evidence that multiple ownership alone makes the system Byzantine.

\subsection{C9. Conflicts and agreements between swarms}
\textbf{Question.} What agreements survive changing membership or access when different swarms have different objectives and can adapt strategically?

\textbf{Obligation.} Choose a concrete peaceful shared resource and one stability concept before comparing mechanisms. State utilities, information, permitted unilateral or coalition deviations and the time during which earlier commitments bind. Prove when a boundary abstraction preserves the deviations that determine the stability claim, or give a parameterized failure and a repaired agreement. A matched pair with the same apparent outcome but different executable deviations is a useful test. Compare with \href{https://www.sciencedirect.com/science/article/pii/S016726810200015X}{dynamic coalition formation}; transferring a basic game formulation to robots is not a research result on its own.

\subsection{C10. Live synthesis while work continues}
\textbf{Question.} In a setting with time to compute, how should a system operate while it constructs or checks a revised policy?

\textbf{Obligation.} Specify what arrives only after the change, a safe interim policy, what work it may continue, and what options it must preserve until the new controller can take over. Characterize an admissible installation condition and a fallback if synthesis times out or the model changes again. Compare a precompiled contingency policy, bounded online replanning, and a genuinely new synthesized policy with identical observations and compute budgets. Find a setting where online computation adds a reachable service option or avoids a genuine precompilation cost, or show that it does not. \href{https://arxiv.org/abs/2107.01136}{Live Synthesis} already handles unfinished obligations at update; the question here includes the behavior during computation and distributed installation. Do not put this branch on the flight-critical path by default.

\subsection{C11. Cascades during a live handoff}
\label{sec:cascade}
\textbf{Question.} When two changes interact during a live handoff, how far must the response reach, and which proofs survive the repair?

\textbf{C11-R: Reachability and interaction.} Freeze a plant, nominal controller, initial envelope, safety invariant and service promise. Specify two primitive changes and their onset rules. From a common reachable prefix with outstanding work, replay the nominal, first-only, second-only and combined continuations under a matched exogenous schedule. The first three must meet the named guarantee on that schedule; the fourth must fail it through traced dependencies across at least two distinct downstream interfaces. For the single-change tolerance claimed beyond the witness, give a separate quantified argument. Exhibit a second reachable cut where the same edits do not require the same downstream intervention. State a nondegenerate parameter region or explain the exact event-order boundary. A supplied local state, individually possible events on different executions, and failure to re-establish a sufficient bound do not witness a cascade.

\textbf{C11-P: Repair and necessity.} On the same model and envelope, construct a policy installed in advance that handles the combined change with the original safety and service promises, under actual local observation and staggered detection. Show a nonempty set of unchanged controllers exposed to the edited subsystem or resource, and prove that at least one specific narrower repair class cannot meet the promise. Explain whether the bound is about observations, authority, capacity or time. A changed future input may alter a fixed controller's behavior; distinguish that from rewriting the controller and from reopening its proof. Prove the repair for the declared history class, not just the witness trace. A well-typed impossibility may close a stated branch but does not deliver a positive repair. The companion specifies the admissible exclusions and positive test.

\textbf{C11-T: Transfer and proof work.} Use a pinned Contextual Composition result in the proof of an unchanged joint guarantee. Build its raw execution, shared boundary, joint law or admissible relation, and required compatibility from the model and source edit; do not supply an already checked whole-system guarantee as a premise. Freeze the interface and proof before an independent selection of a second change. Carry the same proof through that change, or localize the exact premise that fails and demonstrate a checked repair of the interface for a positive transfer. Compare with an already modular direct proof on the same plant, histories and obligations. Report common setup, new local and cross-component proofs, and validation separately. A tie is an admissible result. Apparent savings obtained by moving work into an adapter or weakening the promise do not count.

\subsection{C12. Legitimate changes of goal}
\textbf{Question.} When operators change the work in progress, who can revise its obligations and what must still happen?

\textbf{Obligation.} Specify authorization, effective time, notification and what can be withdrawn, transferred or discharged. Compare two reachable histories with the same new goal and different pending duties. Prove a locally executable transition or characterize the missing authorization or evidence, while preserving the original obligations until a valid revision takes effect. Include one obligation whose physical effect outlasts the authority to start more work.

\subsection{C13. Legitimate changes of safety or risk tolerance}
\textbf{Question.} What transition is required when an authorized decision changes the permissible exposure, even if evidence and the mission do not change?

\textbf{Obligation.} Fix the loss, time horizon, decision authority and moment the new rule binds. For both a tightening and a relaxation, distinguish permissible new actions from exposure already incurred. Prove a transition that preserves the governing bounds over mixed old and new authority, or characterize states from which a newly effective bound cannot be met. Compare the same revision at two reached cuts with different residual exposure. Unlike C6, the requirement itself changes here.

\section{Integrated work and completion}
\textbf{Reuse inside the project.} Use \href{https://github.com/eternal-november/eos-research-programmes/blob/e1348964f73704212d2d8b48d5982a5ed600a2ff/projects/contextual-composition/foundations/shared-world-native-assembly/Whole_Scope_Block_Decomposition.tex#L31-L51}{Whole Scope Block Decomposition} where the shared-resource law yields a native assembled behavior, or \href{https://github.com/eternal-november/eos-research-programmes/blob/e1348964f73704212d2d8b48d5982a5ed600a2ff/projects/contextual-composition/feedback-and-concurrency/observable-covers/Residual_Profiles_and_Admission_Reflection.tex#L84-L101}{Residual Profiles and Admission Reflection} where a changed consumer must use the existing interface. The latter does not prove that newly reached histories satisfy the old laws. An alternative pinned result is admissible if its application carries the stated guarantee and the comparison exhibits what depends on it. A syntax rewrite, a projection of an already solved global law or citation-only reuse is insufficient. The companion records precise source, check and rejection criteria.

\textbf{First integrated return.} Complete C11-R, C11-P and C11-T on one simulation platform. Use two materially different workloads or a held-out change and workload that jointly vary the operational dependence, not just the numeric seed. Use an exact small reference and a separate event or motion-and-message simulator with a published replay. Include work in progress, at least two traced downstream interfaces beyond the directly edited region, and affected and genuinely unchanged parts on the same execution. Report safety, timed service, refused work, extra reserve and verified proof reuse under matched information and resources. No physical trials or operational readiness claims. The outstanding families remain active research questions; completion of this return does not close them.

\textbf{Review and scientific disposition.} Freeze the design envelope, direct baseline and acceptance tests before the main run and independently select the transfer after the interface is fixed. Have a non-author replay the same reference and audit the scope of reuse and the strongest technically compliant shortcut. Report separately whether the undertaking was completed, whether each theorem or obstruction was established, whether an actual proof saving was found, and whether any claim is new relative to a named prior result. A matched direct proof can tie without invalidating the result. An obstruction or prior result can close an investigation at its stated scope without satisfying a required positive deliverable. The companion \emph{Cascade Acceptance Protocol v0.5} defines the first return's submission and review procedure.

\end{document}
